\section{Further research questions}
\label{sec:further}

The new result isolates several precise next problems.
The questions distinguish state reduction, logical
calibration, and proof engineering rather than treating
all three as interchangeable.

\begin{enumerate}[label=\textbf{Q\arabic*.},leftmargin=2.8em]
\item \textbf{Close the exact 238-state proof-assistant endpoint.}
Can the positive-block theorem and its guard be ported
compactly into the current Lean framework, or is direct
new-table coverage cheaper to formalize? The exact
source, initializer, row map, guard events, and four
local paths are now available. The goal is a theorem
whose table constant is definitionally or propositionally
identified with the supplied 476 rows.

\item \textbf{Internalize the transport theorem.}
Can a single arithmetic proof about finitely coded
positive-block simulations transfer an internal
consistency calibration to every later state reduction?
The new step map has values in $\{1,3,5,9\}$ and
$n\le\tau(n)\le9n$. This provides a concrete small
instance for a uniform proof-code construction.
The new package does not settle the repository's
separate internal PA endpoint.

\item \textbf{Search for 237 states using proven compatibility.}
Can another helper role be shared after the two
accepted reductions, with its caller contexts
certified on the resulting source? The successful
238 construction shows why a search over combinations
of prior reductions deserves attention. Every proposed
combination should regenerate its relevant invariants;
disjoint row edits alone should not be an acceptance
test.

\item \textbf{Derive $\rho_3$ inside the exact reduced numerical calculus.}
Theorem~\ref{thm:collectioncong} supplies a short
ordinary first-order argument. Does a uniform numeric
proof scheme exist using only the other seven logical
rows and legal gated bundle instances? A satisfactory
answer must handle variable collisions, freshness
conditions, and the needed instantiations explicitly.
This is the exact obstacle between the relational
theory theorem and a legitimate emitter modification.

\item \textbf{Compile the relational theory $T_-$ completely.}
Can a common-domain quantifier or clause engine for
$T_-$ cost fewer states than the present logical
engine after all substitutions, term coding,
enumeration, and blank initialization are charged?
The right comparison is between complete literal
machines. Deleting five formula-tree nodes from an
axiom display is only an initial design observation.

\item \textbf{Exploit reachable semantic invariants without losing completeness.}
The present runtime covers arbitrary primary values
at every source leaf, a stronger contract than the
actual proof enumerator necessarily needs. Could a
proved phase invariant exclude enough contexts to
permit another fold? Any exclusion must follow
from source semantics and initialization, not from
the absence of a case in a finite trace.

\item \textbf{Unify source grammar and decoder repair search.}
The boundary source decreased its distinct subtree
count while initially disrupting a decoder context.
Can an optimizer minimize total source/control cost
while tracking finite tape-context contracts as part
of the search state? The retained source certificate
and successful repair provide an exact positive
benchmark for such an optimizer.

\item \textbf{Shrink the verification trust surface.}
Can the affine/list case records be exported into a
small independently checked proof language with
explicit literal, unary-induction, and list-induction
constructors? That would preserve the compact
mathematical reasoning while separating discovery
and certificate generation from acceptance more
sharply.
\end{enumerate}

The concrete deliverable of this investigation is the
238-state table with its ordinary mathematical proof,
complete reproduction package, and precise source
calibration. The collection-congruence theorem is an
additional proved theory simplification. The questions
above describe subsequent research and are not premises
left unresolved in the accepted 238-state construction.
